\documentclass[11pt,a4paper]{article}
\usepackage[utf8]{inputenc}
\usepackage[T1]{fontenc}
\usepackage[portuguese]{babel}
\usepackage{xcolor}
\usepackage{hyperref}
\usepackage{geometry}
\usepackage{tcolorbox}

\geometry{margin=2.5cm}

\definecolor{primary}{RGB}{41, 128, 185}
\definecolor{secondary}{RGB}{44, 62, 80}
\definecolor{codebg}{RGB}{240, 240, 240}

\hypersetup{colorlinks=true, linkcolor=primary, urlcolor=primary}

\title{\color{secondary}\Huge\textbf{Exercício 5.7: Faça o deploy para o serverless}}
\author{\color{primary}DevOps com Kubernetes -- 2026}
\date{}

\begin{document}
\maketitle

\section*{\color{primary}Instruções do Exercício}
\begin{tcolorbox}[colback=blue!5!white,colframe=blue!75!black,title=Exercício 5.7: Faça o deploy para o serverless]
Modifique o serviço Ping-pong do seu aplicativo 'Log Output' para utilizar o modelo \textit{serverless} implantando-o como um Knative Service.

- Utilize a nomenclatura de DNS do serviço totalmente qualificado (\textit{fully qualified name}, ex: \texttt{http://pingpong.exercises.svc.cluster.local}) para garantir que os demais componentes locais consigam se comunicar com ele na infraestrutura do Knative.

- Defina o \texttt{hostname} nas regras do \texttt{URLRewrite} caso queira testar a exposição do endpoint \textit{pingpong} e torná-lo acessível no seu navegador.
\end{tcolorbox}

\end{document}
